\section{Discussion}
\label{sec:discussion}

\subsection{Answering the question}

Can chain-of-thought serve as an interpretability layer for Alzheimer's biomarker models? In this
setting, no. CoT did not match black-box accuracy on any of the three tasks, it did not reliably
beat its own no-reasoning controls, and on the primary amyloid/tau task its chains failed the causal
battery decisively. The combination that matters, though, is not ``CoT is inaccurate''---that could
be fixed with a better model. It is \emph{high rated plausibility with near-zero causal
dependence}: 3.67/4 on a rubric with a validated discriminative control, alongside a 3.6\%
sensitivity to a tenfold corruption of the chain's own content. An explanation that is persuasive
and inert is worse than no explanation, because it transfers unearned confidence to the clinician.

\subsection{What survives}

Three components of the hypothesis survive in modified form, and they are why we regard this as a
warning rather than a null result.

\para{Robustness is partly vindicated, but the lesson is about tabular foundation models.} The LLM
arm is genuinely flat under missingness where \xgb collapses, and there is a real crossover on
brain age beyond roughly 25\% missingness. This matches \citet{kearney2026tapgpt}: LLM prompts
represent missing cells natively and need no imputation. But \tabpfn dominates the LLM at every
level tested, and on \taskthree the \cotarm arm is the least robust of all. ``Use an LLM for
robustness'' is not the right conclusion; ``use a tabular foundation model'' is.

\para{\cfit supports a weaker, more interesting claim.} Even when the chain as a whole is post-hoc,
the individual measurements it names are 1.4--2.1$\times$ more influential than the ones it ignores.
Local citations are informative even when the global narrative is not. This suggests a constrained
citation format---``name the three measurements you used'' plus the answer---might retain most of
the interpretability value at a fraction of the accuracy cost of free-form narration.

\para{Calibration improves under CoT on \taskthree} (ECE 0.083 vs.\ 0.314). Communicating
uncertainty well is a genuine clinical good, and it is a benefit the interpretability hypothesis did
not anticipate.

\subsection{The finding we did not anticipate}

CoT faithfulness is markedly higher for regression than for classification (\aoc 0.378 vs.\ 0.066;
mistake sensitivity 0.325 vs.\ 0.036). The entire faithfulness literature is built on
multiple-choice benchmarks, and we had to invent a continuous-target version of the battery to ask
the question at all. If this replicates, it implies that the field's pessimism about CoT
faithfulness may be partly an artefact of its answer format---and that CoT's interpretability value
may be highest precisely on the continuous clinical endpoints (biomarker levels, atrophy rates,
time-to-conversion) that matter most in AD. It is also the finding we consider most likely to
replicate and most worth testing first.

\subsection{Relation to prior work}

We corroborate \citet{anon2026dementiaspeech}, who found text rationales gave inferior
dementia-classification accuracy versus LLM-free baselines, and extend that from speech to
quantitative biomarker tables. We corroborate and sharpen the size-versus-faithfulness result of
\citet{lanham2023faithfulness}, now with the accuracy control \citet{bentham2024disguised} demanded
and in a clinical domain: at 1.5B, CoT is both useful and faithful; at 7B it is neither. And we
supply the measurement that \citet{ding2025l2ctabpfn} and \citet{kearney2026tapgpt} both omit.
\citet{kearney2026tapgpt} disabled reasoning modes; our results suggest that decision was correct
for accuracy purposes, and that had they enabled reasoning, the resulting explanations would have
looked convincing while being causally inert.

Separately, our amendment A2 is a reportable data-provenance finding: an \ltoc-style pipeline
applied naively to OASIS-2 yields AUC $=1.000$ for every model, because OASIS derives its dementia
label from CDR and the lagged CDR features are the label. Anyone replicating
\citet{ding2025l2ctabpfn} on OASIS should de-circularise the feature block before reporting numbers.

\subsection{Limitations}

\para{Data and scope.} ADNI, TADPOLE, AIBL and OASIS-3 are all access-gated, so we replicated
protocols on open cohorts; absolute numbers are \emph{not} comparable to the published tables of
\citet{ding2025l2ctabpfn} or \citet{kearney2026tapgpt}, and only within-study contrasts are valid.
\act is post-mortem tissue rather than in-vivo PET, so \taskone tests neuropathological rather than
imaging \atn, and its age variable is censored above 89. The cohorts are small (107 donors, 150
subjects, 207 subjects), so CIs are wide and a $\Delta$AUC of 0.03 is not resolvable here---which is
why equivalence margins were pre-declared rather than chosen post hoc. We had no plasma biomarkers,
no PET and no multi-omics. We tested one model family at two sizes; a frontier reasoning model might
behave differently, and no API keys were available in our environment, so this is the single most
important follow-up.

\para{Methodology.} The clinician panel could not be convened and was not simulated. The rubric is
scored by an LLM, and by the same base model that produced the chains, which risks self-preference;
the shuffled control establishes that the rubric discriminates, not that its absolute scores match
expert judgement. \cfit and \esa depend on a hand-written lexical citation table that is
deliberately conservative, so both are biased toward the null. The \qwensmall \cotarm arm has a 45.6\%
forced-answer rate, so its faithfulness numbers are computed on chains that were sometimes truncated
before the forced answer; three independent metrics support the direction of the model-size
conclusion, but its magnitude is uncertain. The LLM robustness sweep used a 150-case subset and
fewer perturbation levels than the tabular sweep, so the LLM noise slopes are two-point fits and
should be read as coarse. The brain-age bias correction is fitted on control out-of-fold predictions
and the controls are then one of the two compared groups---standard practice, but mildly circular,
which makes the demented-group gap ($+4.26$\,y for \tabpfn) the more robust statistic than the AUC.
Two targets or feature sets were amended after inspecting baseline behaviour; both changes were made
before any LLM was run and both are audited in the results, but they are researcher degrees of
freedom and we disclose them as such. Finally, greedy decoding with a single run removes sampling as
a confound between arms but gives no seed-to-seed variance estimate; our uncertainty is quantified
over cases, not over generations.

\para{What would invalidate these conclusions.} A frontier reasoning model showing high faithfulness
\emph{and} competitive accuracy on the same three tasks would overturn the central claim; our result
may be specific to a 7B model that was not reasoning-tuned.

\subsection{Implications for clinical deployment}

Any deployment of CoT as a clinical interpretability layer should be required to report a
faithfulness measurement alongside its accuracy. The battery used here costs roughly six extra
generations per case and is fully automatable, so there is no practical excuse for omitting it. On
this evidence, post-hoc SHAP over a tabular foundation model remains both the more accurate and the
more honest option. The broader point generalises past this domain: if a chain can be simultaneously
3.67/4 plausible and 3.6\% causal, then human-rated interpretability and machine faithfulness are
measuring different things, and clinical AI standards that rely on the former are measuring the
wrong one.
